% RESULTADO_RECUPERADO desde II REV09. Véase MAPA_PROCEDENCIA.json.
% Selección de líneas y cambios mecánicos explícitos; ninguna prueba nueva.
\section{Acción, gravedad e información térmica}
\label{v:ant:sec:gravity}
\subsection{Período dodecafásico y realización circular de la acción}
La composición longitudinal de~\eqref{xii:radial:length} determina
\(L=L_*\alpha^{16}5759/23040\), y su reloj compatible satisface
\(t_0=\pi L/(54c_{\rm int})\). El calendario ya construido determina
$T_{108}=108t_0$. En esta sección se transporta esa realización al
lector circular, conservando \(R_{108}=L\).
La velocidad constitutiva, la duración y la longitud de trayectoria se
obtienen mediante las identidades de la
proposición~\ref{iii:prop:identidad-velocidad} y de las
ecuaciones~\eqref{iii:eq:velocidad-longitud-reloj} y
\eqref{iii:eq:longitud-elemental-constitutiva}.
En su realización circular,
\[
\omega_{108}=\frac{2\pi}{108t_0},\qquad
R_{108}=\frac{c_{\rm int}}{\omega_{108}}
=\frac{54}{\pi}\ell_0,\qquad \ell_0=c_{\rm int}t_0.
\]
La unidad $\ell_0$ y las coordenatizaciones de tiempo y velocidad se mantienen
explícitas; la coordenatización SI no selecciona el ciclo.
Para cada sección orientada de acción $a\in\{\mathrm{pre},\mathrm{ret}\}$,
\[
E_{108,a}=\hbar_a\omega_{108},\qquad
m_{108,a}=\frac{E_{108,a}}{c_{\rm int}^2}.
\]
La longitud reducida y la completa son respectivamente
\[
\frac{\hbar_a}{m_{108,a}c_{\rm int}}=R_{108},
\qquad
\frac{h_a}{m_{108,a}c_{\rm int}}=2\pi R_{108}.
\]
La misma acción y la misma duración intervienen en ambas lecturas.

\begin{proposition}[Selector de coincidencia de radios]
En la colección dimensional
\[
G_a(\vartheta)=\vartheta\frac{c_{\rm int}^5t_0^2}{\hbar_a},
\qquad \vartheta>0,
\]
la condición de la realización circular
$G_a(\vartheta)m_{108,a}/c_{\rm int}^2=R_{108}$
tiene la única solución
\[
\vartheta^\circ_{108}=(54/\pi)^2.
\]
\end{proposition}
\begin{proof}
El primer radio es
$r_g(\vartheta)=\vartheta(\pi/54)\ell_0$.
Compararlo con $(54/\pi)\ell_0$, con $\ell_0>0$, reduce la
igualdad a una ecuación lineal estrictamente creciente.
\end{proof}
La coincidencia circular es la restricción al modo unidad del lector
radial positivo del teorema~\ref{xii:radial:rigidity}. Su solución
conserva las composiciones longitudinal y métrica allí declaradas.
El radio reducido es \(Gm/c^2\); el radio de Schwarzschild es
\(2Gm/c^2\) y transmite su correspondiente factor dos.

\begin{proposition}[Acción y gravedad recíprocas, área conservada]
En las dos coordenatizaciones anteriores,
\[
\frac{G_{\rm pre}}{G_{\rm ret}}=R_{\rm act}^{-1},\qquad
\hbar_aG_a=\vartheta^\circ_{108}c_{\rm int}^5t_0^2,
\qquad
\ell_P^2=\frac{\hbar_aG_a}{c_{\rm int}^3}
=\vartheta^\circ_{108}\ell_0^2.
\]
\end{proposition}
\begin{proof}
La definición de $G_a$ contiene $\hbar_a^{-1}$ y el resto de factores
es común a las dos coordenatizaciones. Tomar el cociente y el producto da
las tres igualdades.
\end{proof}
Esta reciprocidad explica por qué la información de orientación puede
variar sin que cambie el área. Se dispone así de una relación precisa
entre las secciones recíprocas de acción, la escala de Planck y el operador areal
de la sección~\ref{v:ant:sec:area}.
